% Unit: Computer Science A --- collections, basic structures, complexity, and patterns (\input after \texttt{cs-a/04-practice.tex} in \texttt{main.tex}).

\runningheader{Computer Science A}{Unit 5: Collections and algorithms}{Page \thepage\ of \numpages}

\begingradingrange{csau5}

\uplevel{\par\textbf{The Java Collections Framework}\par}

\question[4]
What is a \texttt{Collection}?
\makeemptybox{1.15in}
\begin{solution}
  Root interface in the Java Collections Framework for groups of objects (\texttt{add}, \texttt{remove}, \texttt{size}, iteration); \texttt{Map} is not a \texttt{Collection} but is part of the same library.
\end{solution}

\question[4]
What is the Collection API?
\makeemptybox{1.2in}
\begin{solution}
  The set of interfaces and implementations in \texttt{java.util} (and related) for lists, sets, queues, maps, sorting, and utilities such as \texttt{Collections} and \texttt{Arrays}.
\end{solution}

\question[4]
What are the major differences between a \texttt{List}, \texttt{Queue}, \texttt{Set}, and \texttt{Map}?
\makeemptybox{1.65in}
\begin{solution}
  \texttt{List}: ordered, indexed, duplicates allowed. \texttt{Queue}: FIFO (or priority) processing order. \texttt{Set}: no duplicates, often unordered (\texttt{TreeSet} sorted). \texttt{Map}: key--value pairs, keys unique; not a \texttt{Collection} of entries in the same sense.
\end{solution}

\question[4]
What is the difference between a \texttt{Stack} and a \texttt{Queue}?
\makeemptybox{1.2in}
\begin{solution}
  Stack is LIFO (last in, first out); queue is FIFO (first in, first out). Prefer \texttt{Deque} over legacy \texttt{Stack} in modern Java.
\end{solution}

\newpage

\uplevel{\par\textbf{Data structures and algorithms}\par}

\question[4]
Describe a linked list.
\makeemptybox{1.35in}
\begin{solution}
  Nodes where each holds data and a reference to the next (and previous for doubly linked); supports \(O(1)\) insert/delete at a known node but \(O(n)\) indexed access unless optimized.
\end{solution}

\question[4]
What is an algorithm?
\makeemptybox{1.05in}
\begin{solution}
  A finite, well-defined procedure that transforms input into output; analyzable for correctness and resource use.
\end{solution}

\question[4]
How would you describe time complexity?
\makeemptybox{1.35in}
\begin{solution}
  Big-\(O\) (and related) notation for how runtime grows with input size \(n\), ignoring constants and lower-order terms; describes worst/average cases as appropriate.
\end{solution}

\question[4]
What is the difference between linear search and binary search?
\makeemptybox{1.35in}
\begin{solution}
  Linear scans each element in order---\(O(n)\). Binary search halves a sorted range each step---\(O(\log n)\) but requires sorted (or indexed) data.
\end{solution}

\question[4]
What is the difference between Bubble Sort and Merge Sort?
\makeemptybox{1.45in}
\begin{solution}
  Bubble sort repeatedly swaps adjacent out-of-order elements; simple but typically \(O(n^2)\). Merge sort divides and merges sorted halves; \(O(n\log n)\) time, uses extra space for merging.
\end{solution}

\newpage

\uplevel{\par\textbf{Design patterns}\par}

\question[4]
What is the Factory design pattern?
\makeemptybox{1.25in}
\begin{solution}
  Encapsulates object creation behind a method or factory class so callers depend on interfaces/supertypes and creation rules stay in one place (variants: simple factory, factory method, abstract factory).
\end{solution}

\question[4]
What is the Singleton design pattern?
\makeemptybox{1.25in}
\begin{solution}
  Ensures a class has at most one instance and provides global access; useful for shared resources but easy to misuse (testing, threading)---often replaced with DI-managed beans.
\end{solution}

\endgradingrange{csau5}
